\section{Small orbits have sparse supports}\label{sec:support}

The automorphisms used in this section have two parts. Face translations
form an abelian quotient, but their commutators produce central flows on
edges. This interaction constrains the stabilizer of an object with a small
orbit. We will turn that constraint into a short tuple of atoms whose
pointwise stabilizer in the central flow subgroup fixes the object.

Throughout this section, $\mathbb F=\mathbb F_3$, $L\ge2$, and $G=(V,E)$ is
the positively oriented grid on $\{0,\ldots,L-1\}^3$. Its incidence map is
$\partial:\mathbb F^E\to\mathbb F^V$. For a unit square face $f$, write
$c^f\in\mathbb F^E$ for its oriented boundary, so that $\partial c^f=0$.
Let $F$ denote the set of faces and put
\[
 C_F=(\mathbb F^2)^F,\qquad K=\ker\partial.
\]
We use the alternating form
$\omega((a,b),(a',b'))=ab'-ba'$ on $\mathbb F^2$.

\subsection{A central extension acting on the structures}

For $a,a'\in C_F$, define an edge vector
\[
 B(a,a')=\frac12\sum_{f\in F}\omega(a_f,a'_f)c^f\in K.
\]
On the set $H=C_F\times K$, define
\begin{equation}\label{eq:H-product}
 (a,k)(a',k')=(a+a',k+k'+B(a,a')).
\end{equation}
Bilinearity gives
$B(a,a')+B(a+a',a'')=B(a',a'')+B(a,a'+a'')$, proving
associativity. The identity is $(0,0)$; since $B(a,a)=0$, the inverse of
$(a,k)$ is $(-a,-k)$. Thus $H$ is a finite group. The subgroup
$\{0\}\times K$, identified with $K$, is central, and the projection
$H\to C_F$ is a surjective homomorphism with kernel $K$. With the convention
$[h,h']=hh'h^{-1}(h')^{-1}$, its commutator is
\begin{equation}\label{eq:H-commutator}
 [(a,k),(a',k')]=\left(0,\sum_{f\in F}
                  \omega(a_f,a'_f)c^f\right).
\end{equation}

Recall that an edge atom on $e$ has a state $(u,z)$ in $P_e$, with
$u=(u_f)_{f\in F(e)}$. The element $(a,k)\in H$ sends that state to
\begin{equation}\label{eq:H-edge-action}
 \left(u+(a_f)_{f\in F(e)},\;
 z+k_e+\frac12\sum_{f\in F(e)}c^f_e\omega(a_f,u_f)\right).
\end{equation}
It acts on a configuration by applying this map to every incident edge
state. This is an action because restricting $(a,k)$ to an edge is a
homomorphism from $H$ to $P_e$, and the edge action is left multiplication.

\begin{lemma}\label{lem:H-automorphisms}
For every $b\in\mathbb F^V$, the action just defined is by automorphisms of
$A_b$. On its central subgroup $K$, the action adds $k_e$ to each edge
coordinate $z_e$ and leaves all face coordinates unchanged.
\end{lemma}
\begin{proof}
At a vertex $v$, the two incident coordinates belonging to a face $f$
are equal before the action and both receive the same addition $a_f$.
Hence the face compatibility equations remain true. The change in the
divergence of a configuration is
\[
 (\partial k)_v+\frac12\sum_{f\ni v}
   \left(\sum_{e\ni v}\epsilon_{ve}c^f_e\right)
                        \omega(a_f,u_f).
\]
Here $u_f$ is the common face coordinate in that configuration. Both terms
vanish: $k\in K$ and $\partial c^f=0$. Configurations therefore remain in
the same vertex block and still have divergence $b_v$. Applying the inverse
group element proves bijectivity.

The relations $\mathsf{Ed},\mathsf{Cf},\mathsf{EB},\mathsf{VB}$ are
preserved because types and blocks are preserved. The incidence relation
$\mathsf I$ is preserved because configurations and their specified edge
neighbors are transformed by the same map. Finally, for edge states
$g,g'\in P_e$ and $h\in P_e$, one has $(hg)^{-1}(hg')=g^{-1}g'$.
Consequently every $\mathsf Z_\delta$ is preserved. The formula for the
central action follows from \eqref{eq:H-edge-action} with $a=0$.
\end{proof}

Every permutation of the atoms extends to hereditarily finite objects by
$h\cdot x=\{h\cdot y:y\in x\}$ for sets $x$, recursively in their rank.
Atoms remain distinct from sets. We always use this extension of the
$H$-action. A tuple $\alpha$ of atoms is a \emph{$K$-support} of an object
$x$ if every element of the pointwise stabilizer
\[
 K_\alpha=\{k\in K:k\cdot\alpha_i=\alpha_i\text{ for every }i\}
\]
fixes $x$. The empty tuple is permitted and has stabilizer $K$.

\subsection{A sparse representative for an edge functional}

A central element $k\in K$ fixes an edge atom on $e$ exactly when
$k_e=0$. To obtain a short support for an object $x$, it suffices to find
a small edge set $T$ such that every $k\in K$ vanishing on $T$ fixes $x$.
The stabilizer of $x$ in $K$ is a vector subspace, so it can be described
by linear equations on $K$. We will represent these equations using few
edge coordinates.

An edge cochain is a vector $\ell\in\mathbb F^E$. Its face curl is
\[
 (D\ell)(f)=\sum_e c^f_e\ell_e.
\]
A vertex function $\varphi\in\mathbb F^V$ has gradient
$\partial^{\mathsf T}\varphi$, whose value on an edge from $v$ to $w$
is $\varphi(w)-\varphi(v)$. Adding a gradient does not change the curl,
because $\partial c^f=0$. It also does not change the linear functional
$k\mapsto\sum_e\ell_e k_e$ on $K$.
The next lemma shows that a cochain with few nonzero face curls can be
replaced by one using few edges, without changing its functional on $K$.
In the support proof, commutators
will supply the bound on the number of these faces.

\begin{lemma}[Sparse cochains]\label{lem:sparse-cochain}
Let $\ell\in\mathbb F^E$ have nonzero curl on exactly $t$ faces of the
$L\times L\times L$ box. There is a vertex function $\varphi$ such that
$r=\ell-\partial^{\mathsf T}\varphi$ is supported on at most $2Lt$ edges.
In particular, if $t=0$, then $\ell$ is a gradient.
\end{lemma}
\begin{proof}
We construct the subtracted gradient by successive prefix sums.
Write a vertex as $(i,j,k)$. The cochain coordinate $\ell_1(i,j,k)$ is
on the edge from $(i,j,k)$ to $(i+1,j,k)$, and similarly for directions
$2$ and $3$. Define
\[
 \varphi_1(i,j,k)=\sum_{a=0}^{i-1}\ell_1(a,j,k)
\]
with an empty sum interpreted as zero. Subtracting its gradient produces
a cochain $r^{(1)}$ whose direction-$1$ coordinates vanish. Thus
\begin{align*}
 (D\ell)_{12}(i,j,k)
   &=r^{(1)}_2(i+1,j,k)-r^{(1)}_2(i,j,k),\\
 (D\ell)_{13}(i,j,k)
   &=r^{(1)}_3(i+1,j,k)-r^{(1)}_3(i,j,k).
\end{align*}
Here a subscript $ab$ denotes a face in coordinate directions $a<b$,
and the lower corner ranges over the indices for which the face exists.

Let $t_{ab}$ be the number of nonzero $ab$ curls. The $L$ slices of fixed
first coordinate partition the $23$ faces, so choose $i_0$ for which the
number $t_0$ of nonzero $23$ curls satisfies $t_0\le t_{23}/L$. Put
\[
 \varphi_2(i,j,k)=\sum_{b=0}^{j-1}r^{(1)}_2(i_0,b,k),
 \qquad r^{(2)}=r^{(1)}-\partial^{\mathsf T}\varphi_2.
\]
The function $\varphi_2$ is independent of $i$. Consequently $r^{(2)}_1=0$,
and $r^{(2)}_2(i_0,j,k)=0$. Next put
\[
 \varphi_3(i,j,k)=\sum_{c=0}^{k-1}r^{(2)}_3(i_0,0,c),
 \qquad r=r^{(2)}-\partial^{\mathsf T}\varphi_3.
\]
This last function is independent of $i,j$. We have
\[
 r_1=0,\qquad r_2(i_0,j,k)=0,\qquad r_3(i_0,0,k)=0.
\]
On the slice $i=i_0$, the $23$ curl identity now reads
\[
 r_3(i_0,j+1,k)-r_3(i_0,j,k)=(D\ell)_{23}(i_0,j,k).
\]
It follows that a nonzero value of $r_3$ on this slice must lie after a
nonzero curl on its $j$-line. Each nonzero curl accounts for at most $L$
positions. The slice therefore contains at most $Lt_0$ nonzero edge
coordinates; its direction-$2$ coordinates are all zero. Extending these
slice values constantly across all $L$ slices uses at most
$L^2t_0\le Lt_{23}$ supported coordinates.

The difference between $r_2(i,j,k)$ and $r_2(i_0,j,k)$ is a signed sum of
the $12$ curls between $i_0$ and $i$. Each nonzero $12$ curl can occur in
at most $L$ of these sums. Thus at most $Lt_{12}$ direction-$2$ coordinates
differ from their copied slice values. The identical argument in direction
$3$ gives at most $Lt_{13}$ differing coordinates. Since direction $1$
is zero, the support of $r$ has size at most
\[
 L^2t_0+L(t_{12}+t_{13})\le Lt\le2Lt.
\]
The total subtracted function is
$\varphi=\varphi_1+\varphi_2+\varphi_3$. This construction also covers
$t=0$, when the resulting support is empty.
\end{proof}

\subsection{The support bound}

\begin{proposition}[Small-orbit support bound]\label{prop:small-orbit-support}
Let $N=|A_b|$, let $q\ge0$, and let $x$ be a hereditarily finite object
over the atoms of $A_b$. If its $H$-orbit has size at most $N^q$, then
$x$ has a $K$-support consisting of at most
\[
 2L(q\log_3 N)^2
\]
edge atoms. The support may be empty.
\end{proposition}
\begin{proof}
Let $S=\{h\in H:h\cdot x=x\}$, let $U\le C_F$ be the image of $S$
under the projection, and put $K_0=S\cap K$. Since $C_F$ and $K$ are
additive groups over the prime field $\mathbb F_3$, both $U$ and $K_0$
are vector subspaces. We use the group exact sequence with kernel $K_0$ to obtain
$|S|=|U|\,|K_0|$. Therefore, with
$a=\operatorname{codim}_{C_F}U$ and $d=\operatorname{codim}_K K_0$,
\begin{equation}\label{eq:stabilizer-codimension}
 a+d=\log_3[H:S]\le q\log_3 N=:h.
\end{equation}

Let $W\le K^*$ be the annihilator of $K_0$, so $\dim W=d$. Every
$\lambda\in K^*$ extends to a functional on $\mathbb F^E$, represented
by an edge cochain $\ell$. Two such cochains differ by an element of
$K^\perp=\operatorname{im}\partial^{\mathsf T}$. Indeed that image is
contained in $K^\perp$, and both have dimension $\operatorname{rank}\partial$.

Fix $\lambda\in W$ and a representative $\ell$. For $u,u'\in U$,
choose lifts in $S$. Their commutator lies in $S\cap K=K_0$; evaluating
$\lambda$ on \eqref{eq:H-commutator} gives
\[
 \sum_{f\in F}(D\ell)(f)\,\omega(u_f,u'_f)=0.
\]
Thus $U$ is totally isotropic for the alternating form on $C_F$ with
one block $(D\ell)(f)\omega$ for each face. If $t$ of these coefficients
are nonzero, the form has rank $2t$: a nonzero scalar multiple of
$\omega$ has rank two, and the blocks have disjoint face coordinates
in $C_F$. This rank calculation does not require the boundary vectors
$c^f$ in $K$ to be linearly independent.

We recall the standard dimension argument for alternating forms
\cite[Sections~3.1 and~3.3]{cameron2000classical}.
An alternating form of rank $2t$ on an
$n$-dimensional vector space has no totally isotropic subspace of
codimension less than $t$. Quotient by the radical, of dimension
$n-2t$. The image of an isotropic subspace in this nondegenerate
$2t$-dimensional space has dimension at most $t$, since it is contained
in its orthogonal complement. The original isotropic subspace therefore
has dimension at most $(n-2t)+t=n-t$. Applied above, this shows $t\le a$.
By Lemma~\ref{lem:sparse-cochain}, $\lambda$ consequently has a cochain
representative supported on at most $2La$ edges.

Choose such representatives for a basis of $W$, and let $T\subseteq E$
be the union of their supports. Then
\[
 |T|\le2Lad\le2Lh^2.
\]
This includes $d=0$, when the basis and $T$ are empty. Any $k\in K$
that vanishes on $T$ annihilates every basis member of $W$, hence all
of $W$. Finite-dimensional duality gives $k\in K_0$, so $k$ fixes $x$.
For each edge $e\in T$, choose one edge atom. By the central action,
fixing that atom is equivalent to $k_e=0$. The tuple of chosen atoms
is therefore a $K$-support of $x$ and has the stated length.
\end{proof}

The argument used $x$ only through its stabilizer subgroup $S\le H$.
Its rank plays no role in the support bound.

A family $\mathcal X$ is \emph{membership-transitive} if
$y\in x\in\mathcal X$ implies $y\in\mathcal X$. Thus it contains
every recursive constituent of each of its members.

\begin{corollary}\label{cor:transitive-support}
Let $\mathcal X$ be an $H$-invariant, membership-transitive finite family
of hereditarily finite objects over $A_b$, and suppose
$|\mathcal X|\le N^q$ for some fixed $q\ge0$. Every object in
$\mathcal X$, including every recursive constituent of each of its
members, has a $K$-support of length at most $2L(q\log_3N)^2$.
All these supports can be padded to the common positive length
\begin{equation}\label{eq:support-length}
 s=\left\lceil2L(q\log_3N+1)^2\right\rceil+1
   =O_q\bigl(L(1+\log L)^2\bigr).
\end{equation}
\end{corollary}
\begin{proof}
The orbit of each member is contained in $\mathcal X$, so the proposition
applies to it. Membership transitivity places all its recursive
constituents in the same family. Adding atom coordinates to a support
can only shrink its pointwise stabilizer, so it remains a support.
There are atoms because $L\ge2$; arbitrary atoms can therefore pad even
an empty support. Finally $N=\Theta(L^3)$ gives the asymptotic estimate.
\end{proof}
